\documentclass[conference]{IEEEtran}
\usepackage{cite}
\usepackage{amsmath,amssymb,amsfonts}
\usepackage{booktabs}
\usepackage{array}
\usepackage{multirow}
\usepackage{url}
\usepackage{xcolor}
\newcommand{\indicator}{\mathbb{I}}
\begin{document}

\title{Why Do Humans Stop Growing Taller?\\
A Longitudinal Proposal for Growth-Plate Closure Readiness}

\author{\IEEEauthorblockN{Anonymous Research Plan Author}\\
\IEEEauthorblockA{Proposal generated from a frozen Survey--Idea--ExperimentDesign record\\
Design-only scientific research plan}}

\maketitle

\begin{abstract}
The statement that genes determine when people stop growing is an incomplete account of a developmental process. This proposal operationalizes the question as the transition from measurable longitudinal long-bone growth to a sustained low-growth state. The biological end organ is the epiphyseal growth plate, where chondrocyte proliferation and differentiation are regulated by local pathways and systemic developmental signals. Human estrogen-resistance evidence and experimental growth-plate studies support an important role for pubertal estrogen signaling, while genetic studies support a polygenic contribution to adult stature rather than a deterministic individual closure clock. We propose a noninterventional longitudinal Growth-Plate Closure Readiness Atlas. Repeated stadiometry, radiation-free magnetic-resonance imaging features, puberty-related measurements, and separately consented inherited feature summaries will be evaluated against a chronological-age-plus-bone-age baseline. The primary endpoint is a preregistered sustained low-growth state; the secondary structural endpoint is interval-censored growth-plate closure. Locked temporal and external-site validation, feature-block ablations, subgroup calibration gates, and strict pediatric governance define what would count as evidence. No participant data, fitted model, causal estimate, or clinical recommendation is reported. The contribution is a falsifiable protocol for separating plausible biological mechanisms from validated individual-level prediction.
\end{abstract}

\begin{IEEEkeywords}
growth plate, epiphyseal fusion, longitudinal bone growth, puberty, estrogen, skeletal maturation, predictive modeling, calibration
\end{IEEEkeywords}

\section{Introduction}

People commonly use ``stop growing'' to mean that a person stops getting taller. That language hides several distinct biological processes. Body mass can change throughout adulthood, bone remodels continuously, and tissues can respond to injury or disease. None of those processes is the endpoint studied here. This proposal concerns the termination of longitudinal growth in long bones: the transition from measurable height gain to a sustained low-growth state after the developmental activity of the epiphyseal growth plate has diminished and, in many sites, fusion has occurred.

The simplified explanation that genes determine the stopping point contains a grain of truth but is scientifically insufficient. Inherited variation contributes to body size and to pathways that influence developmental timing. However, it is neither a complete individual clock nor a substitute for observing growth velocity, puberty, skeletal state, health context, and measurement quality. Large-scale human genetic evidence describes highly polygenic adult-height architecture rather than a single closure gene \cite{wood2014}. A scientific account must also accommodate the central role of growth-plate biology and pubertal endocrine signaling.

The growth plate is a layered cartilaginous structure at the ends of growing bones. Resting-zone cells contribute progenitors; proliferative chondrocytes create longitudinal columns; hypertrophic chondrocytes differentiate and are replaced by bone. Local pathways, including PTHrP and Indian hedgehog (IHH), coordinate this process with endocrine and growth-factor signaling \cite{vanderEerden2003,kronenberg2006}. Thus, chronological age is a proxy for a changing biological state, not the state itself.

Estrogen signaling provides a particularly informative example. A man with a disruptive estrogen-receptor mutation showed incomplete epiphyseal closure and continued linear growth into adulthood, illustrating that estrogen is important for skeletal maturation in humans of more than one sex \cite{smith1994}. Experimental work showed that estrogen accelerated multiple growth-plate senescence markers and was associated with epiphyseal fusion \cite{weise2001}. Related rabbit evidence proposed irreversible depletion of resting-zone progenitors as a mechanism by which estrogen hastens fusion \cite{nilsson2014}. The human clinical observation and animal mechanisms are mutually informative, but they must not be collapsed into an unmeasured universal dose-response curve for children.

This paper presents a design-only research proposal rather than an executed study. Its central question is narrower than ``why do we stop growing?'' in a philosophical sense: \emph{can a preregistered, longitudinal, noninterventional model of growth-plate and pubertal state improve calibrated prediction of sustained longitudinal growth cessation beyond chronological age and bone age alone?} The answer will remain conditional on a prospectively validated protocol. The proposal does not predict a particular participant's adult height, recommend hormones, or make a clinical decision.

\section{Background, Survey, and Research Gap}

\subsection{Growth plate as the terminal organ of longitudinal growth}

Longitudinal bone growth depends on growth-plate cell kinetics and architecture. The plate integrates systemic signals such as growth hormone and IGF-related inputs with local regulation by PTHrP, IHH, fibroblast growth factors, bone morphogenetic proteins, and extracellular-matrix processes \cite{vanderEerden2003,kronenberg2006}. These pathways affect the pace at which chondrocytes proliferate, differentiate, hypertrophy, and are replaced by bone. A model that uses only age can therefore be useful as a population proxy but is unable to observe the end organ whose state gives age its predictive value.

Growth-plate senescence is a developmental decline in proliferative capacity. It is not simply the passing of calendar time. Nilsson and Baron summarized evidence that plate chondrocytes have finite proliferative capacity and that local maturation helps explain changes in catch-up growth and eventual fusion \cite{nilsson2005}. This distinction matters for study design: two adolescents of the same chronological age can have different pubertal timing, height velocity, skeletal maturity, and remaining growth potential. It also matters for interpretation: an apparent age association cannot identify a molecular or endocrine cause without a design that measures the relevant state variables.

\subsection{Estrogen, sex, and species boundaries}

Estrogen is a key driver of skeletal maturation, but an accurate account must avoid two common errors. First, estrogen biology is not restricted to females; the human estrogen-receptor case supports a role in male skeletal maturation as well \cite{smith1994}. Second, evidence that a manipulation changes fusion in an animal model does not establish a human intervention effect. The PNAS experiment showed accelerated senescent decline and fusion after estrogen treatment \cite{weise2001}; the rabbit study connected this acceleration to a loss of resting-zone progenitors \cite{nilsson2014}. These sources motivate candidate measurements and a mechanistic diagram, not a prescription for manipulating puberty.

The proposed protocol therefore records puberty-related endocrine measurements only as potential predictive covariates. It does not infer that a serum value alone determines closure, because hormone exposure at the growth plate depends on timing, tissue response, receptor function, feedback loops, and other biology. It does not require a research biopsy from healthy minors. The design uses repeated, radiation-free imaging and standard anthropometry to test whether the integrated state adds predictive information.

\subsection{Genes influence the system without defining a complete clock}

Adult height has a broad polygenic architecture, and associated loci map to diverse biological functions \cite{wood2014}. This supports an inherited contribution to the geometry and regulation of growth, but it does not make a polygenic score a valid ``time-to-stop-growing'' device. Effects can be small, population-dependent, correlated with ancestry or environment, and unstable across sites. A predictive study must therefore treat inherited summaries as separately consented, quality-controlled candidate inputs and audit their incremental value and subgroup calibration. A model that gains average discrimination while worsening reliability for a subgroup is not acceptable as the primary model.

\subsection{Research gap}

The Survey identifies two actionable gaps. \textbf{GAP-GP-01} is a missing longitudinal human alignment of repeated height velocity, noninvasive growth-plate state, puberty-related measurements, and inherited regulatory features prior to sustained cessation. Much available evidence is mechanistic, cross-sectional, or based on exceptional cases. \textbf{GAP-GP-02} is the absence of a preregistered, out-of-sample-calibrated protocol that tests whether an integrated closure-readiness model improves on chronological age and bone age. A one-time bone-age association does not resolve either gap because it does not measure future residual growth or report probability calibration.

\begin{table*}[t]
\caption{Frozen Survey evidence and the permitted design use}
\label{tab:evidence}
\centering
\footnotesize
\begin{tabular}{p{0.18\textwidth}p{0.31\textwidth}p{0.42\textwidth}}
\toprule
Source & Verified scope & Permitted use in this proposal \\
\midrule
Smith \cite{smith1994} & Human estrogen-receptor disruption with incomplete epiphyseal closure and continued linear growth. & Motivate estrogen-related puberty features in all sexes; no therapeutic inference. \\
Weise \cite{weise2001} & Experimental estrogen accelerated growth-plate senescence and fusion. & Motivate a maturation mechanism; no human effect-size claim. \\
Nilsson \cite{nilsson2014} & Rabbit evidence connecting estrogen exposure, earlier fusion, and resting-zone progenitor depletion. & Species-bounded mechanistic rationale only. \\
van der Eerden and Kronenberg \cite{vanderEerden2003,kronenberg2006} & Local and systemic growth-plate regulation, including PTHrP and IHH. & Define multi-block feature families and biological interpretation. \\
Nilsson and Baron \cite{nilsson2005} & Local growth-plate senescence and finite proliferative capacity. & Reject a universal chronological-age endpoint. \\
Wood \cite{wood2014} & Polygenic adult-height architecture. & Restrict genetics to an ancestry-aware optional predictive block. \\
\bottomrule
\end{tabular}
\end{table*}

\section{Research Questions and Planned Contributions}

The primary research question is: \emph{after all preprocessing choices are locked, does a longitudinal multimodal model improve held-out calibration and proper scoring for sustained low longitudinal growth beyond a chronological-age-plus-bone-age baseline?} Four supporting questions make this testable.

\begin{enumerate}
\item Can repeated MRI growth-plate features and puberty-related measurements improve primary held-out probability calibration beyond the baseline?
\item Does any incremental value remain stable across temporal and external-site holdouts?
\item Are predicted probabilities calibrated across prespecified developmental-stage, sex-variable, site, and ancestry strata with sufficient data support?
\item Does an optional inherited-feature block add stable predictive value after quality control and ancestry-aware controls, or should it be removed for parsimony?
\end{enumerate}

The proposal makes five planned contributions. First, it separates the narrow endpoint of longitudinal growth cessation from broader changes in adult bodies. Second, it keeps mechanistic evidence and predictive claims distinct. Third, it defines a common baseline and a set of feature-block ablations so a complex model cannot claim success through a changed outcome or data split. Fourth, it incorporates interval-censored structural closure and sustained height-velocity information rather than using a single cross-sectional label. Fifth, it makes subgroup calibration, consent, data minimization, and specialist review formal release gates rather than afterthoughts.

\section{Idea Source Checkpoints and Direction Selection Audit}

The Idea stage started only after the Survey registry was frozen. Three routes were compared. The first was a cross-sectional skeletal-maturity catalogue. It would be inexpensive and visually attractive, but a single age, bone-age, or imaging measurement can only describe present association; it cannot test whether a marker predicts future residual growth. It was not selected.

The selected route is the \emph{Longitudinal Growth-Plate Closure Readiness Atlas}. It follows repeated standardized stadiometry, MRI features, puberty-related measures, and limited covariates through a developmental window. It estimates a probability of later sustained low growth and an interval-censored closure time, then tests whether each additional feature block improves held-out performance. This directly addresses both accepted gaps while avoiding intervention.

The third route was direct endocrine perturbation and tissue sampling. Such a study might yield mechanistic information, but elective hormone manipulation or growth-plate biopsy is not an acceptable initial method for healthy minors. It was rejected for ethics, feasibility, and interpretability. A high-risk ``universal genetic closure clock'' was also rejected because the Survey supports polygenic architecture, not deterministic personal closure timing.

No Monte Carlo tree search, evolutionary optimization, or autonomous idea evolution was run. The route-selection record is a transparent, evidence-constrained comparison. This is important because inventing an optimization history would create false provenance without adding scientific value.

\section{Problem Definition, Assumptions, and Hypotheses}

\subsection{Operational endpoints}

Let $H_i(t)$ denote the calibrated standing height of participant $i$ at visit time $t$. For adjacent scheduled visits, the observed velocity is
\begin{equation}
V_i(t)=\frac{H_i(t)-H_i(t-\Delta t)}{\Delta t},
\label{eq:velocity}
\end{equation}
where $\Delta t$ is the registered interval in years. Height collection follows a single stadiometer protocol and site-quality procedure. The primary endpoint is a sustained low-growth indicator,
\begin{equation}
Y_i(t)=\indicator\{V_i(t)<v^* \ \mathrm{and}\ V_i(t+\Delta t)<v^*\},
\label{eq:primary_endpoint}
\end{equation}
where $v^*$ is a threshold selected before outcome inspection. It is a research definition used to compare models, not a universal biological or clinical cutoff.

For the secondary structural endpoint, the actual closure time $T_i$ is not observed continuously. If MRI indicates an open plate at visit $L_i$ and closed plate at the next adequate scan $R_i$, then
\begin{equation}
T_i \in (L_i,R_i].
\label{eq:interval_censoring}
\end{equation}
The interval-censoring representation prevents an artificial claim that closure occurred on the imaging date. Participants with persistent open plates at the study horizon are right-censored, not assigned an invented closure time.

At each landmark time $t$, the candidate model produces
\begin{equation}
p_i(t)=\Pr\left[Y_i(t+\Delta t)=1 \mid \mathcal{H}_i(t),\mathcal{M}\right],
\label{eq:prediction}
\end{equation}
where $\mathcal{H}_i(t)$ contains only data available at or before $t$ and $\mathcal{M}$ is the frozen model specification. The baseline contains chronological age, standardized height history, recorded sex variables, and protocol-approved bone age. The multimodal candidate adds MRI features and puberty-related measurements; inherited summaries, if available under separate consent, are a final ablation rather than a default input.

\subsection{Assumptions and falsifiable hypotheses}

The study assumes that height measurements are standardized, MRI quality is recorded, visit timing is captured, and any bone-age assessment is acquired under an approved clinical or research protocol. It does not assume that every growth plate is visible with equal quality, that every participant provides blood or genetic data, or that a single biological pathway has the same predictive relation in all populations. Missingness, consent choices, site effects, and subgroup support are part of the analysis rather than nuisances to be hidden.

\textbf{H1:} The MRI-plus-puberty feature blocks improve primary held-out Brier score and calibration over the baseline. \textbf{H2:} The selected model remains acceptably calibrated across registered sites and supported subgroups. \textbf{H3:} The inherited-feature block is retained only if it adds repeatable value after ancestry-aware controls. Each can fail. A null incremental benefit, an external-site failure, or a subgroup calibration defect is an informative result that prevents a stronger predictive claim.

\begin{table}[t]
\caption{Registered variables and roles}
\label{tab:variables}
\centering
\footnotesize
\begin{tabular}{p{0.28\columnwidth}p{0.62\columnwidth}}
\toprule
Element & Operational role \\
\midrule
Height history $H_i(t)$ & Standardized repeated anthropometry used for velocity and residual-growth endpoints. \\
MRI feature vector $X_i^{MRI}(t)$ & Prespecified noninvasive growth-plate morphology or signal features after quality control. \\
Puberty vector $X_i^{pub}(t)$ & Protocol-approved puberty stage and endocrine measurements; predictive, not causal labels. \\
Bone age $B_i(t)$ & Baseline proxy when available under the approved protocol. \\
Inherited summary $G_i$ & Separately consented, ancestry-aware optional feature block. \\
Closure interval $(L_i,R_i]$ & Interval-censored structural outcome derived from adequate serial MRI. \\
\bottomrule
\end{tabular}
\end{table}

\section{Study Design and Methods}

\subsection{Cohort, consent, and acquisition}

The intended design is a prospective, noninterventional cohort spanning preregistered prepubertal through late-pubertal developmental windows. The unit of analysis is a participant-visit, while validation splits preserve participant-level dependence. The final enrollment target is not guessed in this proposal. Before recruitment, a simulation will specify the number of participants and visits needed for event-rate uncertainty, site clustering, attrition, missingness, and subgroup calibration precision. A study that cannot support its declared subgroup analyses will narrow those claims rather than report unstable estimates.

At each registered visit, staff measure standing height using calibrated equipment and a fixed protocol. A radiation-free MRI acquisition targets a prespecified long-bone growth plate; image quality, sequence parameters, and reasons for unusable scans are logged. Puberty-related measurements are collected only under pediatric ethics approval and with participant assent and guardian consent as applicable. Samples and image data follow separate minimization and access controls. Genetic collection has separate opt-in consent, defined retention, a withdrawal process, and no use outside the registered purpose.

Figure \ref{fig:pipeline} summarizes the proposed information flow. It is a protocol diagram, not a result figure. The design excludes research biopsy, experimental hormone administration, and use of model outputs to recommend care. Any incidental finding follows the approved clinical pathway and is never managed by the research model.

\begin{figure*}[t]
\centering
\fbox{\begin{minipage}{0.94\textwidth}
\centering\small
\textbf{Prospective visits} $\rightarrow$ standardized height and visit timing $\rightarrow$ MRI quality-controlled growth-plate features\\[3pt]
$\downarrow$ \hspace{3.5cm} $\downarrow$ \hspace{3.5cm} $\downarrow$\\[-1pt]
protocol-approved puberty measures \hspace{1cm} optional separately consented inherited summaries \hspace{1cm} protected data manifest\\[4pt]
$\rightarrow$ locked baseline and multimodal models $\rightarrow$ temporal and external-site holdouts $\rightarrow$ calibration, subgroup, and human-review gates
\end{minipage}}
\caption{Conceptual pipeline for the proposed Growth-Plate Closure Readiness Atlas. The arrows denote data provenance and validation order, not causal effects.}
\label{fig:pipeline}
\end{figure*}

\subsection{Preprocessing and model comparison}

All feature definitions are versioned before the first holdout is examined. Height records are screened with documented device and protocol checks; no outlier is removed without a recorded reason. MRI preprocessing is performed without access to future outcomes, and image-quality exclusions are applied identically across model variants. The data manifest records enrollment site, visit date, imaging protocol, consent tier, feature availability, missingness reason, transformation version, and linkage key. This provenance is necessary because a nominally accurate model can be driven by site-specific acquisition artifacts.

The baseline model receives only its registered age-and-bone-age inputs and is fitted using the same training partition as every competitor. The candidate model has the same endpoint and data cutoff but adds feature blocks in a locked order. One implementation may use a joint longitudinal and time-to-event formulation in which residual growth and structural closure share participant-level latent effects. A generic closure hazard can be written as
\begin{equation}
\lambda_i(t)=\lambda_0(t)\exp\left\{\beta^\top Z_i(t)+b_i\right\},
\label{eq:hazard}
\end{equation}
where $\lambda_0(t)$ is a baseline hazard, $Z_i(t)$ contains only registered covariates observed by time $t$, $\beta$ is estimated from training data, and $b_i$ is a participant-level random effect. Equation \eqref{eq:hazard} describes a predictive association, not the biological causal mechanism of closure.

The model portfolio contains: (i) age-plus-bone-age baseline, (ii) baseline plus MRI, (iii) baseline plus puberty-related measurements, (iv) full multimodal model, and (v) full model without inherited summaries. This ablation design establishes whether any apparent gain arises from a particular block rather than from a changed cohort, endpoint, or tuning budget. Complex models cannot select their own validation period, imaging features, or decision threshold after observing held-out outcomes.

\subsection{Validation, calibration, and decision gates}

The data are partitioned into a locked training set, a later temporal holdout, and an external-site holdout. Feature selection, imputation models, transformations, tuning, and threshold choice are learned or frozen using training data only. Repeated visits from an individual remain on one side of every split. If an external site uses an incompatible MRI acquisition or cannot provide a required outcome, that limitation is reported rather than silently pooled away.

For a held-out binary endpoint with model probability $p_i$ and outcome $y_i$, the Brier score is
\begin{equation}
\mathrm{BS}=\frac{1}{n}\sum_{i=1}^{n}(p_i-y_i)^2.
\label{eq:brier}
\end{equation}
Here, lower values are better only when compared over the same locked outcomes and partition. The score is reported with calibration intercept and slope, reliability plots, log score, and uncertainty intervals. For residual growth quantities, interval coverage and proper scoring complement the binary endpoint. Overall discrimination alone cannot satisfy the primary decision gate.

Let $K$ indicate that all calibration requirements are met, $E$ indicate that external-site performance meets a registered transportability margin, and $F$ indicate that subgroup audits find no unresolved reliability failure. Let $\delta^*$ be the preregistered practical improvement margin relative to the baseline. The model-release decision is
\begin{equation}
d=\begin{cases}
\text{multimodal}, & \Delta\mathrm{BS}\leq-\delta^*,\ K=E=F=1,\\
\text{baseline}, & \Delta\mathrm{BS}>-\delta^*,\ K=E=F=1,\\
\text{inconclusive}, & \text{otherwise},
\end{cases}
\label{eq:decision}
\end{equation}
where $\Delta\mathrm{BS}=\mathrm{BS}_{\mathrm{multi}}-\mathrm{BS}_{\mathrm{base}}$. Equation \eqref{eq:decision} constrains what may be claimed; it does not determine which result will occur.

\begin{table*}[t]
\caption{Execution stages, controls, and stopping conditions}
\label{tab:protocol}
\centering
\footnotesize
\begin{tabular}{p{0.16\textwidth}p{0.30\textwidth}p{0.27\textwidth}p{0.17\textwidth}}
\toprule
Stage & Planned action & Controls and diagnostics & Stop or fail condition \\
\midrule
1. Governance and registration & Obtain ethics approvals; freeze population, endpoints, MRI protocol, consent tiers, and analysis plan. & Assent/consent review, data minimization, versioned protocol. & Unresolved pediatric ethics, privacy, or clinical oversight issue. \\
2. Acquisition and manifest & Collect serial height, MRI, and approved covariates; document quality and missingness. & Calibration records, image QC, provenance ledger. & Systematic protocol drift or unusable endpoint ascertainment. \\
3. Lock validation & Define training, temporal, and external-site partitions before fitting. & Participant-aware split check and leakage audit. & Split contamination or post-hoc threshold change. \\
4. Fit and ablate & Fit baseline and feature-block models using common endpoints and tuning budget. & Convergence, missing-data sensitivity, predeclared feature order. & Nonconvergence, instability, or data leakage. \\
5. Evaluate & Score calibration, Brier and log scores, coverage, and subgroup reliability. & Site and developmental-stage audit with uncertainty. & Calibration, transportability, or subgroup gate fails. \\
6. Interpret & Apply Eq.~\eqref{eq:decision}; archive provenance and review limitations. & Specialist review and communication check. & Any unresolved ethical or inferential boundary. \\
\bottomrule
\end{tabular}
\end{table*}

\subsection{Sensitivity, missingness, and reproducibility}

The primary analysis does not silently convert unavailable measurements into normal values. A data-availability matrix will distinguish missed visits, image-quality exclusion, declined blood collection, optional genetic nonconsent, laboratory failure, and incompatible acquisition. The complete-case analysis is reported alongside a registered missing-data sensitivity analysis. Any imputation model is fitted inside the training partition and applied without refitting to holdouts. A conclusion that depends on one imputation method or on excluding a site is classified as unstable.

Several further sensitivity analyses are locked before execution. The first varies only the research low-growth threshold $v^*$ across a clinically reviewed, preregistered range and checks whether the model-ranking decision changes. The second changes the landmark and follow-up spacing only when the corresponding visits satisfy the same measurement-quality criteria. The third repeats the analysis with and without participants whose bone-age assessment is unavailable, so a baseline advantage cannot arise solely from selective access to imaging. The fourth tests acquisition-site harmonization by fitting site-aware nuisance adjustments in training data and then evaluating a held-out site without recalibration. These analyses are not opportunities to search for a preferred result; they are diagnostic conditions for narrowing a claim.

Reproducibility requires more than storing model code. Before fitting, the team archives the data dictionary, protocol version, consent-tier map, data schema, split identifiers, preprocessing specification, MRI feature extractor version, random seeds, software environment, and analysis report template. A release record links every reported table and plot to a locked configuration. Human-readable audit notes explain any data-quality exclusion. Deidentified analytic data are shared only when consent, governance, and re-identification risk permit; otherwise a controlled-access or synthetic-data validation route is documented. These steps make it possible to distinguish biological nonreplication from a changed endpoint, image pipeline, or validation split.

\subsection{Interpretation safeguards}

The project registers a strict separation between three levels of statement. A \emph{measurement statement} reports how reliably an MRI or endocrine feature was acquired. A \emph{predictive statement} reports held-out probability performance conditional on the enrolled cohort and protocol. A \emph{mechanistic statement} requires evidence beyond predictive importance and must retain the human-versus-animal boundary in the Survey. Feature importance is therefore not translated into a claim that a variable biologically commands closure. Similarly, a calibrated probability is not a forecast of personal destiny and must never be shown to families as a medical conclusion without an approved clinical framework.

\section{Expected Outcome Branches and Conditional Conclusions}

The study has three predeclared outcome branches. In the first, the multimodal model improves proper scoring and remains calibrated in temporal and external-site validation without unresolved subgroup defects. The appropriate conclusion would be limited: the registered measurements provide conditional predictive information beyond age and bone age in the studied settings. It would not establish that a particular hormone, MRI feature, or genotype causes a participant's growth plate to close.

In the second branch, the age-and-bone-age baseline performs as well as the multimodal model. This would favor parsimony and identify the additional acquisition burden as unjustified for the defined endpoint. It would not refute the biology of growth plates or estrogen; it would only show that the added measurements did not produce validated predictive improvement under this protocol.

In the third branch, neither model meets calibration or transportability gates. This is a scientifically useful outcome because it exposes inadequate endpoint measurement, site heterogeneity, missingness, or an overambitious prediction target. The correct conclusion would be that the present design cannot support reliable individual-level closure-readiness probabilities. No branch allows a personal treatment recommendation, a deterministic genetic label, or an assertion that all humans stop growing by one chronological age.

\section{Risks, Limitations, and Human Review Requirements}

The principal scientific risk is conflating predictive association with causal explanation. Estrogen-related measurements may correlate with closure readiness without identifying the pathway or timing of effect in a person. Animal evidence concerning resting-zone progenitors must remain species-bounded. Genetic features may introduce population-structure artifacts and may be interpreted by families as destiny. The analysis will therefore report feature families, uncertainty, calibration, and failure cases without converting a score into a statement about a child's potential or worth.

The principal measurement risk is endpoint and acquisition quality. Standing-height velocity is sensitive to protocol error and visit spacing. MRI visibility and sequence consistency can vary by site and developmental stage. Structural closure is interval-censored, not directly timed. Missing blood draws or optional genetic consent can be informative and can distort comparisons. These risks are addressed through an explicit data manifest, quality-control logs, sensitivity analyses, and a rule that unsupported subgroup claims are withheld.

The principal ethical risk concerns research with minors. Participation requires age-appropriate assent, guardian consent, privacy protection, and a clinical pathway for incidental findings. MRI avoids ionizing radiation but still demands comfort, safety, and quality procedures. No research biopsies, no experimental endocrine exposures, and no automated clinical actions are permitted. Pediatric endocrinologists, musculoskeletal radiologists, statisticians, privacy specialists, and ethics reviewers must approve execution and interpretation. Until that review and the subsequent study occur, this manuscript remains a proposal with no observed results.

\begin{table}[t]
\caption{Inference boundaries and required review}
\label{tab:risks}
\centering
\footnotesize
\begin{tabular}{p{0.28\columnwidth}p{0.62\columnwidth}}
\toprule
Risk & Required safeguard \\
\midrule
Age treated as biology & Compare age proxy with direct longitudinal measurements and report incremental value. \\
Animal mechanism treated as human effect & Mark species boundary; avoid human dose or intervention claims. \\
Genetic determinism or population bias & Separate consent, ancestry-aware analysis, subgroup calibration, and human review. \\
Imaging or visit artifacts & Lock MRI QC, document missingness, and use external-site validation. \\
Clinical overreach & Research-only outputs; no automated care or treatment inference. \\
\bottomrule
\end{tabular}
\end{table}

\begin{thebibliography}{00}
\bibitem{smith1994} E. P. Smith \textit{et al.}, ``Estrogen resistance caused by a mutation in the estrogen-receptor gene in a man,'' \textit{New England Journal of Medicine}, vol. 331, no. 16, pp. 1056--1061, 1994, doi: 10.1056/NEJM199410203311604.
\bibitem{weise2001} M. Weise, S. De-Levi, K. M. Barnes, R. I. Gafni, V. Abad, and J. Baron, ``Effects of estrogen on growth plate senescence and epiphyseal fusion,'' \textit{Proceedings of the National Academy of Sciences}, vol. 98, no. 12, pp. 6871--6876, 2001, doi: 10.1073/pnas.121180498.
\bibitem{nilsson2014} O. Nilsson \textit{et al.}, ``Evidence that estrogen hastens epiphyseal fusion and cessation of longitudinal bone growth by irreversibly depleting the number of resting zone progenitor cells in female rabbits,'' \textit{Endocrinology}, vol. 155, no. 8, pp. 2892--2899, 2014, doi: 10.1210/en.2013-2175.
\bibitem{vanderEerden2003} B. C. J. van der Eerden, M. Karperien, and J. M. Wit, ``Systemic and local regulation of the growth plate,'' \textit{Endocrine Reviews}, vol. 24, no. 6, pp. 782--801, 2003, doi: 10.1210/er.2002-0033.
\bibitem{nilsson2005} O. Nilsson and J. Baron, ``Impact of growth plate senescence on catch-up growth and epiphyseal fusion,'' \textit{Pediatric Nephrology}, vol. 20, no. 3, pp. 319--322, 2005, doi: 10.1007/s00467-004-1689-4.
\bibitem{wood2014} A. R. Wood \textit{et al.}, ``Defining the role of common variation in the genomic and biological architecture of adult human height,'' \textit{Nature Genetics}, vol. 46, no. 11, pp. 1173--1186, 2014, doi: 10.1038/ng.3097.
\bibitem{kronenberg2006} H. M. Kronenberg, ``PTHrP and skeletal development,'' \textit{Annals of the New York Academy of Sciences}, vol. 1068, pp. 1--13, 2006, doi: 10.1196/annals.1346.002.
\end{thebibliography}

\appendices
\section{Idea-Evolution Record}

The Idea stage compared a cross-sectional catalogue, the selected longitudinal readiness atlas, and direct perturbation with tissue sampling. The catalogue was rejected because present association does not validate future prediction. The perturbation route was rejected because it would expose healthy minors to unacceptable risk and would not be an appropriate initial study. The selected route is noninterventional and remains explicitly predictive. No autonomous search trajectory was executed or claimed.

\section{Symbols and Operational Definitions}

\begin{table}[h]
\caption{Symbols used in the proposed analysis}
\label{tab:symbols}
\centering
\footnotesize
\begin{tabular}{p{0.28\columnwidth}p{0.62\columnwidth}}
\toprule
Symbol & Definition \\
\midrule
$H_i(t)$ & Standardized standing height for participant $i$ at visit time $t$. \\
$V_i(t)$ & Registered interval height velocity from Eq.~\eqref{eq:velocity}. \\
$Y_i(t)$ & Sustained low-growth endpoint from Eq.~\eqref{eq:primary_endpoint}. \\
$T_i$ & Unobserved structural closure time; interval-censored by serial MRI. \\
$p_i(t)$ & Probability of later sustained low growth given information available at $t$. \\
$Z_i(t)$ & Frozen time-indexed predictive covariate vector. \\
$\lambda_i(t)$ & Modelled closure hazard used for predictive association. \\
$K,E,F$ & Calibration, external transportability, and subgroup fairness release gates. \\
\bottomrule
\end{tabular}
\end{table}

\section{Evidence Coverage and Review Checklist}

The Author stage used only the seven citation keys frozen by Survey: smith1994, weise2001, nilsson2014, vanderEerden2003, nilsson2005, wood2014, and kronenberg2006. The manuscript introduces no new source, observed result, or empirical performance value. Human clinical evidence supports the importance of estrogen signaling; experimental estrogen and progenitor evidence is marked as nonhuman; reviews define growth-plate regulation and senescence; and the genetic study is limited to polygenic adult-height architecture.

Before execution, reviewers must finalize eligibility, visit spacing, the research endpoint threshold $v^*$, MRI feature extraction, sample-size simulation, subgroup minimum support, external-site comparability, consent language, genetic data governance, and procedures for incidental findings. These reviews are essential because the proposal defines a research protocol, not a clinical product.

\end{document}
